\section{The formal problem: what an API can and cannot see}
\label{sec:setup}
\subsection{Configuration, channels, and the identified set}
Think of an endpoint as a bundle of four things, written $\omega=(M,S,I,P)$: the
\emph{model} itself, the \emph{serving stack} that runs it, the \emph{interface} that
exposes it, and the \emph{preprocessing} that prepares the input.  When we send a
request $X_t$, the API returns $Y_t=(y_t,\ell_t,u_t,t_t,h_t)$.  Those five parts are
the sampled text, the numeric values (logprobs or probabilities), the usage accounting,
the timing, and the envelope (headers and error schemas).  These are our only windows
into the system.

Our target is not $\omega$ itself.  It is the \emph{equivalence class}
$[\omega]_{\mathcal C,\mathcal X}$: the set of every configuration that would be
impossible to tell apart from $\omega$ when viewed through the channels $\mathcal C$ we
have, using only the requests $\mathcal X$ we are allowed to send.  For each property
$g$ we also track the \emph{identified set} $\mathcal I_g(\omega;\mathcal C,\mathcal
X)$: the values of $g$ that remain possible (Def.~\ref{def:identified} in
Sec.~\ref{sec:ident}).  What we recover is \emph{functional} architecture: how the
input is tokenized, how text is generated, whether the model uses attention or a
constant-size state, whether the feed-forward layers are dense or sparse, the geometry of
the output, the numeric regime, and the serving structure.  In other words, the
organisation of the computation, as far as the equivalence class pins it down.  Never
weights, and never source code.  Throughout the paper we call this \emph{functional
architecture identification under API observability constraints}.  When we say
``reconstruction'', that is the shorthand.

\begin{table}[!htb]\centering\footnotesize
\caption{Plain-word glossary.  The terms and grades used throughout the paper, in the
order a reader meets them.  Grades L1 to L5 and U are attached to every reconstructed
component; the four boundary classes I, C, A, N are attached to every open question.}
\label{tab:glossary}
\begin{tabularx}{\textwidth}{@{}l Y@{}}
\toprule
\textbf{Term} & \textbf{Meaning in plain words}\\
\midrule
$\omega=(M,S,I,P)$ & The hidden system: model, serving stack, interface, preprocessing.\\
channel & One kind of thing the API returns that we can measure: usage counts, logprobs, timing, behaviour, error text.\\
$[\omega]_{\mathcal C,\mathcal X}$ & All systems that would look identical to $\omega$ through our channels and allowed requests.  This is what we can actually recover.\\
identified set $\mathcal I_g$ & For one property $g$, the values that are still possible after everything we observed.\\
witness & A concrete pair of different systems that give identical observations.  A witness proves a property cannot be identified.\\
$\kappa$, $\hat\kappa$ & Information per request: how much evidence one API call gives for telling two candidates apart.  High (around 2) means a few requests settle it.  Low (around 0.02) means hundreds.  Zero means never.\\
probe (P01 to P23) & A controlled request pattern designed to measure one thing, with a stated failure mode.\\
support score & The engine's ranking score for a candidate.  It is a ranking, not a probability; we measure how far from a probability it is.\\
prefill / decode & Prefill is the work done on the prompt before the first output token.  Decode is the per-token work after that.\\
KV cache / constant state & A transformer stores a growing key-value cache, so each new token costs a little more.  An SSM keeps a fixed-size state, so it does not.\\
\midrule
L1 observed & A value the API returned directly or that we measured directly (an error message, a token count).\\
L2 identified & A controlled experiment distinguished the stated alternatives.\\
L3 estimated & A numeric estimate with a model and an uncertainty or resolution statement.\\
L4 mechanism-consistent & The evidence favours one candidate, but alternatives remain.\\
U unresolved & The procedure, sample, channel, or budget was not enough to decide.\\
L5 non-identifiable & A scoped witness proves no admissible request can decide.\\
\midrule
I / C / A / N & Boundary classes: Identified; Cost-limited (separable, priced by $\hat\kappa$); Access-limited (the API hides the channel); Non-identifiable (witness).\\
\bottomrule
\end{tabularx}
\end{table}

\subsection{The Forensic Layer Model}
\label{sec:layers}
The single most common mistake in informal probing is \emph{layer confusion}: reading an
interface fact as if it were a model fact.  To guard against this we treat the four
layers as separate causal zones, and every observation is tagged with the layer it is
allowed to speak about:
\begin{center}
\emph{model} $\neq$ \emph{serving} $\neq$ \emph{interface} $\neq$ \emph{preprocessing}.
\end{center}
\begin{itemize}
\item A changed \textbf{output cap} is an interface change, not a model change.  We
  demonstrate this with the cap fixture (Sec.~\ref{sec:blindfixtures}).
\item A change in \textbf{latency} says something about serving, and nothing about the
  model, until queueing, batching, and caching are controlled.  The warm-versus-cold
  contrast in Fig.~\ref{fig:attncontrol} is a factor of 5 to 10 on its own.
\item \textbf{Token accounting} is a statement about the provider's ledger, not the
  tokenizer, until it is checked against a tokenizer we count locally.  The
  ledger-inflation fixture shows why.
\item \textbf{Logprob grids} are a statement about the wire format.  The dtype used for
  the actual computation sits behind the rounding applied in transport.  The dtype axis
  of \benchname{} measures this gap.
\end{itemize}
A provider controls its own ledger \cite{chen2026tokeninflation}.  A client's clock
includes network noise.  A deployment cap (16{,}393 tokens on GPT-4o-mini) is not the
model's capability \cite{openaimini2024docs}.  This layer discipline is what lets partial
recovery stay honest.  Each node of the hypothesis graph carries its provenance
(provider-reported, payload-observed, client-measured, or independently verified) next
to its evidence grade.

\subsection{Access regimes are first-class}
\label{sec:regimes}
Every result is reported together with the access regime it used.  \emph{Standard}
access means text, usage counts, top-$k$ logprobs, and timing, which is what a typical
production API exposes.  \emph{Restricted} access means top-5 logprobs only and no logit
bias.  \emph{Rich} access adds a diagnostic selected-column channel that we built into
our own serving wrapper, purely to measure what wider access would buy.  The regimes are
never mixed.  Production claims use only standard access.  \benchname{} reports coverage
for each regime separately (Sec.~\ref{sec:geometry}).
